%% ma: L12-B08-0026
%% lop: 12
%% chuong: 2
%% bai: 8
%% dang: 12.8.18
%% loai: DS
%% muc: [TH, TH, VD, VD]
%% dap_an: "ĐSĐS"
%% nguon: "(NBV)-ĐỀ SỐ 17-MỖI NGÀY 1 ĐỀ THI 2025.pdf (D17, MỖI NGÀY 1 ĐỀ - Đề số 17), Phần II Câu 3"
%% kien_thuc: "Bài 7 (hệ trục tọa độ), Bài 8 (tọa độ điểm, độ dài vectơ, tích vô hướng); thể tích khối chóp (Lớp 11 Bài 27); góc nhị diện (Lớp 11 Bài 25)"
%% trang_thai: cho-duyet
%% ly_do: "Dạng mới đề xuất 12.8.18 chờ giáo viên duyệt"
%% ghi_chu: "Đáp án gốc ĐSĐS đúng (giải độc lập, kiểm bằng sympy/numpy: SI = √2, V = √2/3, khối lượng ≈ 16,499 g, góc nhị diện = arccos(-1/3) ≈ 109,47°). Hình dựng lại bằng phép chiếu song song với đúng tọa độ các đỉnh; trục Ox hướng về phía người xem; nét khuất (IA, ID, DA) vẽ đứt; các đoạn OC, OD, OS nằm trong khối vẽ đứt. Bỏ ảnh trang trí."
\begin{ex}
Một vật trang trí nhỏ bằng pha lê có dạng hình bát diện đều, mỗi cạnh có độ dài $1$ cm (tức là hình gồm hai hình chóp tứ giác đều chung đáy $S.ABCD$ và $I.ABCD$ gắn với nhau như hình dưới đây, độ dài tất cả các cạnh bên và cạnh đáy đều bằng nhau và bằng $1$ cm). Người ta gắn hệ trục tọa độ $Oxyz$ như hình dưới đây (đơn vị tính theo cm). Biết tọa độ của $C\left(\dfrac{\sqrt2}2;0;0\right)$.
\begin{center}
\begin{tikzpicture}[x={(.6cm,-.5cm)},y={(1cm,.15cm)},z={(0cm,1cm)},scale=3.4]
  \pgfmathsetmacro{\k}{sqrt(2)/2}
  \coordinate (O) at (0,0,0);
  \coordinate (A) at (-\k,0,0);\coordinate (C) at (\k,0,0);
  \coordinate (B) at (0,-\k,0);\coordinate (D) at (0,\k,0);
  \coordinate (S) at (0,0,\k);\coordinate (I) at (0,0,-\k);
  \draw (S)--(A) (S)--(B) (S)--(C) (S)--(D) (B)--(C)--(D) (A)--(B) (I)--(B) (I)--(C);
  \draw[khuat] (D)--(A) (I)--(A) (I)--(D) (O)--(C) (O)--(D) (O)--(S);
  \draw[truc] (C)--(1.05,0,0) node[below,font=\small]{$x$};
  \draw[truc] (D)--(0,1.0,0) node[right,font=\small]{$y$};
  \draw[truc] (S)--(0,0,.95) node[above,font=\small]{$z$};
  \foreach \p/\d in {S/above left,I/below,A/above left,B/left,C/right,D/above right}
    \fill (\p) circle (.3pt) node[\d,font=\footnotesize] {$\p$};
  \fill (O) circle (.3pt) node[below left,font=\footnotesize] {$O$};
\end{tikzpicture}
\end{center}
\choiceTF
{\True Tọa độ điểm $A$ là $\left(-\dfrac{\sqrt2}2;0;0\right)$.}
{Khoảng cách giữa hai đỉnh $S$ và $I$ của khối pha lê là $2$ cm.}
{\True Khối lượng của khối pha lê xấp xỉ $16{,}5\ \mathrm g$ (biết khối lượng riêng của loại pha lê làm vật trang trí đã cho là $35\ \mathrm{g/cm^3}$).}
{Số đo góc nhị diện $[S,CD,I]$, làm tròn đến hàng đơn vị bằng $150^\circ$.}
\loigiai{$O$ là tâm hình vuông $ABCD$ có cạnh $1$, nên $OA=OB=OC=OD=\dfrac{\sqrt2}2$ và $SO=IO=\sqrt{SA^2-OA^2}=\dfrac{\sqrt2}2$. Do đó $A\left(-\dfrac{\sqrt2}2;0;0\right)$, $B\left(0;-\dfrac{\sqrt2}2;0\right)$, $D\left(0;\dfrac{\sqrt2}2;0\right)$, $S\left(0;0;\dfrac{\sqrt2}2\right)$, $I\left(0;0;-\dfrac{\sqrt2}2\right)$.

a) Đúng.

b) $SI=2\cdot\dfrac{\sqrt2}2=\sqrt2\ne2$. Sai.

c) $V=2\cdot\dfrac13\cdot S_{ABCD}\cdot SO=2\cdot\dfrac13\cdot1\cdot\dfrac{\sqrt2}2=\dfrac{\sqrt2}3\ \mathrm{cm^3}$. Khối lượng $35\cdot\dfrac{\sqrt2}3\approx16{,}5\ \mathrm g$. Đúng.

d) Gọi $N$ là trung điểm $CD$; $SN\perp CD$, $IN\perp CD$ nên góc nhị diện $[S,CD,I]$ bằng $\widehat{SNI}$. Với $N\left(\dfrac{\sqrt2}4;\dfrac{\sqrt2}4;0\right)$: $\overrightarrow{NS}=\left(-\dfrac{\sqrt2}4;-\dfrac{\sqrt2}4;\dfrac{\sqrt2}2\right)$, $\overrightarrow{NI}=\left(-\dfrac{\sqrt2}4;-\dfrac{\sqrt2}4;-\dfrac{\sqrt2}2\right)$, $\cos\widehat{SNI}=\dfrac{\frac18+\frac18-\frac12}{\frac34}=-\dfrac13$. Vậy $\widehat{SNI}\approx109^\circ\ne150^\circ$. Sai.}
\end{ex}
